\section{REAL Bench: no reportar solo la puntuación total, sino evaluar la distribución de tareas}

Las páginas web en línea cambian continuamente; si se prueba un agent directamente en sitios reales, aparecen actualizaciones de páginas, cambios de inventario, estados de inicio de sesión y riesgos de seguridad, de modo que los distintos modelos no se comparan en el mismo entorno. El aporte principal de REAL es construir réplicas de sitios de alta fidelidad pero deterministas, para que las trayectorias de acciones puedan ejecutarse de nuevo, los fallos puedan analizarse a posteriori y los modelos puedan compararse sobre el mismo conjunto de tareas.

\lecturefigure{slide-15-real-bench-home.jpg}{REAL: entornos web deterministas de alta fidelidad orientados a Agents}{00:14:43}

\begin{importantbox}{Definición: un entorno determinista no es lo mismo que un entorno simple}
\term{Deterministic simulation} significa que el mismo estado inicial y la misma secuencia de acciones producen resultados reproducibles; aun así, puede incluir páginas complejas, varios sitios y tareas largas. El determinismo resuelve sobre todo la comparabilidad de los experimentos; no garantiza que el entorno cubra todo el ruido del mundo real.
\end{importantbox}

\subsection{De la puntuación total de un modelo a un vector de capacidades por dominio}

La diapositiva del gráfico de anillo muestra la REAL score global de un modelo; el gráfico de radar y el gráfico de barras horizontales revelan las diferencias entre sitios. Al leer la figura, el primer paso no es preguntar «quién quedó primero», sino cómo se construyó el muestreo de tareas, cuántas preguntas tiene cada sitio, si los fallos se concentran en cierto tipo de interacción y si la puntuación global está dominada por tareas sencillas.

Estas tres figuras forman una cadena de diagnóstico que va de lo general a lo particular. La puntuación total sirve para responder «¿el sistema mejoró en conjunto?»; el gráfico de radar por sitio, para responder «¿dónde ocurrió la mejora?», y el gráfico de barras por tarea, para localizar «qué página, herramienta o acción sigue fallando». Si un modelo rinde muy bien en el sitio de compras pero falla de forma persistente en los sitios de calendario y de vuelos, el promedio puede parecer aceptable, pero el equipo de producto debería revisar primero la interpretación de fechas, el estado de identidad, los formularios largos y el cambio entre sitios, en lugar de seguir ampliando el modelo a ciegas.

\lecturefigure{slide-16-real-bench-model-score.jpg}{Puntuación total de los modelos en REAL y estadísticas de tareas por sitio}{00:15:22}

\lecturefigure{slide-17-real-bench-site-radar.jpg}{Perfil de capacidades de distintos Agents en cada sitio}{00:16:16}

\lecturefigure{slide-18-real-bench-task-breakdown.jpg}{Desempeño de los modelos desglosado por sitio web y por tarea}{00:17:22}

Desde la perspectiva de producto, el desglose por sitio y por tarea también determina el orden de las correcciones. Si varios modelos fallan en el mismo sitio web, el problema puede venir del diseño del entorno, de la interfaz de las herramientas o de la ambigüedad de la tarea; si solo un modelo falla en las tareas de formularios largos en todos los sitios, lo más probable es que ese modelo tenga deficiencias en la conservación del estado o en la planificación de acciones. Por eso, una plataforma de evaluación no solo debe producir una tabla de clasificación para la investigación, sino también agrupamientos de fallos accionables, traces reproducibles y contraejemplos mínimos, para que los ingenieros puedan reproducir el problema en lugar de adivinarlo a partir de un video.

Sea $\mathcal{W}$ el conjunto de sitios web, $\mathcal{T}_w$ el conjunto de tareas del sitio $w$ y $s_{w,t}\in\{0,1\}$ el indicador de éxito de una tarea individual. El promedio macro es
\[
S_{\text{macro}}=\frac{1}{|\mathcal{W}|}\sum_{w\in\mathcal{W}}
\frac{1}{|\mathcal{T}_w|}\sum_{t\in\mathcal{T}_w}s_{w,t}.
\]
$S_{\text{macro}}$ da el mismo peso a cada sitio; el promedio micro, en cambio, da más peso a los sitios con más tareas. El reporte debe indicar qué agregación se usa; de lo contrario, el mismo resultado puede contarse como historias distintas.

Además, se dividen las tareas en slices $g\in\mathcal{G}$ según riesgo, longitud y tipo de interacción:
\[
S_g=\mathbb{E}[s\mid g],\qquad
S_{\min}=\min_{g\in\mathcal{G}}S_g.
\]
$S_{\min}$ es el peor slice. Un producto de alto riesgo no puede optimizar solo la media; si las tareas de pago, permisos o eliminación forman el peor grupo, un 90\% global tampoco demuestra que el sistema pueda desplegarse.

\begin{knowledgebox}{Lectura de la figura: qué responde cada una de las tres páginas de resultados}
La página de puntuación total responde «¿en qué nivel se encuentra este sistema en conjunto?»; el gráfico de radar responde «¿las capacidades están equilibradas?», y el gráfico de barras por tarea responde «¿dónde se concentran los fallos?». Las tres deben usarse en conjunto: la puntuación total para ordenar, la distribución para diagnosticar y el peor slice como umbral para la puesta en producción.
\end{knowledgebox}

\begin{warningbox}{Benchmark saturation no equivale a confiabilidad en el mundo abierto}
Las réplicas deterministas permiten pruebas de regresión estables, pero los sitios web reales también incluyen actualizaciones de interfaz, latencia, permisos, personalización, medidas contra la automatización y consecuencias irreversibles. El benchmark debe servir como plataforma de entrenamiento y depuración, no como una forma de convertir directamente la puntuación del leaderboard en una tasa de incidentes reales.
\end{warningbox}

\teachervoice{00:14:03--00:17:58, el ponente subraya que REAL usa réplicas de sitios modernas, complejas y deterministas, y que compara los modelos sitio por sitio. El docente enfatiza: lo esencial de un agent web no es «saber hacer clic en botones», sino poder completar tareas de manera estable en distintos sitios, tareas y estados de error.}

\subsection{Resumen de la sección}

REAL Bench resuelve la reproducibilidad y la visibilidad de la distribución en la evaluación de agents. Un reporte confiable incluye, como mínimo, la puntuación global, el promedio macro por sitio, los slices de tareas, el peor grupo y las causas de fallo a nivel de trayectoria. La siguiente sección aborda AgentQ: cómo lograr que un agent no solo sea evaluado, sino que también aprenda de las trayectorias fallidas.
